\documentclass[a4paper,11pt]{ltjsarticle}
\newcommand{\bookletname}{採点基準}
% LuaLaTeX / UTF-8
\usepackage{amsmath,amssymb,bm}
\usepackage[margin=23mm,top=22mm,bottom=22mm]{geometry}
\usepackage{enumitem,booktabs,array}
\usepackage{fancyhdr}
\usepackage[hidelinks,unicode]{hyperref}
\setlength{\parindent}{1\zw}
\setlength{\parskip}{0.35em}
\setlength{\headheight}{17pt}
\pagestyle{fancy}
\fancyhf{}
\fancyhead[L]{\small 第1回模試\quad 統計応用（医薬生物学）}
\fancyhead[R]{\small\bookletname}
\fancyfoot[C]{\small\thepage}
\renewcommand{\headrulewidth}{0.35pt}
\setlist[enumerate,1]{label=〔\arabic*〕,leftmargin=3.5em,labelsep=0.5em,itemsep=0.85em,topsep=0.7em}
\newcommand{\E}{\operatorname{E}}
\newcommand{\Var}{\operatorname{Var}}
\newcommand{\Cov}{\operatorname{Cov}}
\newcommand{\Prob}{\operatorname{P}}
\newcommand{\qhead}[1]{\section*{問#1}}
\newcommand{\ans}[1]{\subsection*{〔#1〕}}
\newcommand{\points}[1]{\hfill\textnormal{（#1点）}}
\newcommand{\continued}[1]{\noindent\textbf{問#1（続き）}\par\medskip}
\hypersetup{pdfauthor={Codex},pdfsubject={統計検定1級対策・自主学習用模擬試験}}

\setlength{\parskip}{0.2em}
\setlist[itemize]{itemsep=0.15em,parsep=0pt,topsep=0.3em,leftmargin=1.5em}
\renewcommand{\ans}[1]{\par\medskip\noindent\textbf{〔#1〕}\quad}
\renewcommand{\points}[1]{\hfill\textnormal{（#1点）}\par\nobreak}
\hypersetup{pdftitle={第1回模試 統計応用（医薬生物学） 採点基準}}
\begin{document}
\begin{center}
{\Large\bfseries 第1回模試\quad 統計応用（医薬生物学）}\par\medskip
{\large 採点基準・部分点}
\end{center}
\section*{採点の共通方針}
各大問30点，選択した3問の合計90点とする。100点換算値は
\[
\text{換算点}=\frac{\text{選択した3問の得点の合計}}{90}\times100
\]
とする。公式の配点・合否判定を再現するものではない。
\begin{enumerate}[label=\arabic*.,leftmargin=2em,itemsep=1em]
\item 模範解答と異なる方法でも，数学的に正しく，設問で指定された方法に対応していれば同じ点を与える。結論が一致していても，必要な導出や理由がない場合は該当する結論点だけとする。
\item 前の小問の計算誤りをそのまま用いた後続答案では，正しい方法による部分を評価し，同じ数値誤りを繰り返し減点しない。ただし，別の概念上の誤りは別途扱う。
\item 四則演算・平方根・付表で再現できる計算を基準とする。付表の補間や途中の丸めの差だけで減点しない。数値が多少異なっても，検定の判断に影響せず方法が正しければ評価する。
\item 「有意差なし」と「同じであることを示した」，条件付きの指標と集団全体の指標，モデルを仮定した区間と仮定の検証を区別する。
\item 採点単位の中で一部だけ正しい場合は，0.5点刻みの部分点を付けてよい。各小問・大問の上限を超えない。
\end{enumerate}
\section*{配点一覧}
\begin{center}
\begin{tabular}{crrrrrr}
\toprule
大問 & 〔1〕 & 〔2〕 & 〔3〕 & 〔4〕 & 〔5〕 & 合計\\
\midrule
問1 & 7 & 11 & 7 & 5 & -- & 30\\
問2 & 4 & 6 & 10 & 6 & 4 & 30\\
問3 & 4 & 8 & 10 & 8 & -- & 30\\
問4 & 4 & 7 & 5 & 7 & 7 & 30\\
問5 & 8 & 5 & 4 & 5 & 8 & 30\\
\bottomrule
\end{tabular}
\end{center}
\newpage

\qhead{1}
\ans{1}\points{7}
\begin{itemize}
\item リスク集合の対 $(n_0,n_1)$ が時点順に $(3,3),(2,3),(2,3),(2,2)$：各1点，計4点。
\item KM積の構成を示す：1点。
\item $\hat S_0(5)=1/3$，$\hat S_1(5)=4/9$：各1点，計2点。
\end{itemize}
打ち切り・新規登録だけで生存曲線を下げた場合は，KM積の構成点を与えない。

\ans{2}\points{11}
\begin{itemize}
\item リスク集合内の死亡者の条件付き寄与を，ハザードの比として説明：2点。
\item 部分尤度 $\theta^2/\{6(1+\theta)^2(2+3\theta)^2\}$：3点。母数に依存しない定数は省略可。
\item 対数尤度と微分またはスコア方程式：3点。
\item $\hat\theta=\sqrt{2/3}$：1点。
\item 導関数の符号変化，または対数部分尤度の狭義凹性による一意の最大の確認：2点。
\end{itemize}
条件付き寄与の説明には死亡時間の密度全体までの記載を必須としない。
微分が0であることだけでは，最大の確認の2点は与えない。

\ans{3}\points{7}
\begin{itemize}
\item $U(0)=-1/5$：2点。対照群側のスコアなら符号が逆でも整合していれば可。
\item $I(0)=49/50$：2点。
\item 統計量 $2/49$ と，付表3の3.84との比較による非棄却：2点。
\item ログランク・スコアが観測数と期待数の差の和であり，$U(0)$ と一致すると説明：1点。
\end{itemize}
正規近似で $|U|/\sqrt I$ と1.96を比較する方法も同点とする。
非棄却から効果がないと断定した場合は，検定判断の部分を1点までとする。

\ans{4}\points{5}
\begin{itemize}
\item 正しい併合KM値0.400：1点。
\item 誤ったリスク人数 $(8,7,5,4)$ とKM値0.450：2点。
\item 登録まで生存していた条件で選ばれた人を，登録前にもリスク集合に入れる誤りを説明：2点。
\end{itemize}
「分母が違う」とだけ述べた場合，理由の2点のうち1点まで。
\newpage

\qhead{2}
\ans{1}\points{4}
粗死亡率A 0.016，B 0.040：各1点。SMR A 1，B $20/17$：各1点。
率を1000人年当たり16，40と示しても，単位が明確なら同点。

\ans{2}\points{6}
\begin{itemize}
\item 両地域の直接標準化死亡率0.028：2点。
\item $w_k=Y_kr_k/\sum_\ell Y_\ell r_\ell$ を示す：2点。
\item 年齢別比が共通に $(0.8,1.2)$ で，重みがA $(1/2,1/2)$，B $(1/17,16/17)$ と異なることからSMR差を説明：2点。
\end{itemize}
単に「高年層が多いから」だけの場合，最後の項目は1点まで。

\ans{3}\points{10}
\begin{itemize}
\item ポアソン確率関数から独立性を用いて尤度を構成：2点。
\item スコア方程式と最大の確認：2点。死亡総数0の境界を含めて記載すれば満点。正の総数だけ正しい場合は1.5点。
\item $\hat\theta=O/E$ と $\E[\hat\theta]=\theta,\ \Var(\hat\theta)=\theta/E$：2点。
\item デルタ法による対数の漸近分散 $1/(\theta E)$，推定標準誤差 $1/\sqrt O$：2点。
\item 対数尺度の区間を指数変換し，$[0.822,1.217]$ を得る：2点。
\end{itemize}
最後の2点は，区間の式1点，付表を用いた数値1点とする。
0.196を0.20に丸めた $[0.819,1.221]$ も，丸めを明記していれば数値点を与える。
元の尺度で対称な区間だけを答えた場合，指定の対数尺度の区間の点は与えない。

\ans{4}\points{6}
\begin{itemize}
\item 当てはめ値 $(50,50)$：1点。
\item 統計量 $X^2=4$：2点。
\item 自由度1と，2平均に対して1母数を推定したという理由：2点。
\item 付表3の3.84との比較で共通倍率モデルを棄却：1点。
\end{itemize}
自由度だけで理由がない場合，該当2点のうち1点。

\ans{5}\points{4}
\begin{itemize}
\item SMRは平均であり，各年齢層の比が1とは限らず，〔4〕もモデルに疑問を示すと説明：2点。
\item 両地域の年齢別率の推定値は等しく，観測されたSMRの差は重みの違いで説明できると述べる：2点。
\end{itemize}
正誤の結論だけなら各主張0.5点。〔3〕の区間に1が含まれることを，同率の証明として用いない。
\newpage

\qhead{3}
\ans{1}\points{4}
\begin{itemize}
\item 感度 $16/17$：1点。
\item 特異度 $19/43$：1点。
\item 陽性構成割合 $2/3$ と母集団0.20からの偏りを指摘：2点。
\end{itemize}
正しい母集団の指標を答える小問ではなく，選択された300人の単純集計を問う小問である。

\ans{2}\points{8}
\begin{itemize}
\item 陰性者への重み8と陽性者への重み1：2点。
\item 推定表の4セルが陽性行 $(80,120)$，陰性行 $(40,760)$：2点。
\item 感度 $2/3$，特異度 $19/22$，陽性的中率0.4，陰性的中率0.95：各1点，計4点。
\end{itemize}
無作為抽出された陰性者を8倍する方法と，条件付き確率と既知の $q$ から再構成する方法は同点。

\ans{3}\points{10}
\begin{itemize}
\item $p=qa+(1-q)b$：1点。
\item $\hat p=qX_1/200+(1-q)X_0/100$ と不偏性：2点。
\item 独立性から分散 $q^2a(1-a)/200+(1-q)^2b(1-b)/100$ を導出：3点。
\item $\hat p=0.12$ と標準誤差 $\sqrt{0.000352}\simeq0.01876$：2点。
\item 1.96を用いた95\%区間 $[0.083,0.157]$：2点。
\end{itemize}
重みの2乗を落とした場合，分散の導出は1点まで。
1000人すべての疾患が観測済みだとみなして $p(1-p)/1000$ を使った場合，分散の点は与えない。
設問は二項モデルを明記しているため，有限母集団修正は不要。

\ans{4}\points{8}
\begin{itemize}
\item ベイズの公式と条件付き独立性から $v_1h/\{v_1h+v_0(1-h)\}$ を導出：4点。式だけなら2点。
\item $\Prob(D=1\mid S=1,R=1)=a$ を示し，陽性的中率は変わらないと結論：2点。
\item 陰性者内の選択が疾患と関連する場合は条件付き独立性が崩れ，重み8の一律補正を正当化できないと説明：2点。
\end{itemize}
「偏る可能性がある」とだけ述べ，どの仮定が崩れるか説明しない場合，最後の項目は1点まで。
\newpage

\qhead{4}
\ans{1}\points{4}
周辺分散 $\sigma_B^2+\sigma_e^2$：1点。同一施設内の共分散 $\sigma_B^2$：1点。
相関係数 $\rho=\sigma_B^2/(\sigma_B^2+\sigma_e^2)$ と共通の施設効果による説明：2点。

\ans{2}\points{7}
\begin{itemize}
\item $\E[\hat\tau]=\tau$：1点。
\item 施設平均の分散 $\sigma_B^2+\sigma_e^2/m$：2点。
\item $\Var(\hat\tau)=2(\sigma_B^2+\sigma_e^2/m)/K$：2点。
\item $V_{\mathrm{ind}}$ を示し，分散比 $1+(m-1)\rho$ を得る：2点。
\end{itemize}
全測定値の共分散を直接足し合わせた導出も同点。

\ans{3}\points{5}
\begin{itemize}
\item 分散 $19/15$，標準誤差 $\sqrt{19/15}$：1点。
\item $Z\simeq2.666$ と1.96との比較による棄却：2点。
\item 95\%信頼区間 $[0.794,5.206]$：2点。
\end{itemize}
施設効果の母分散も既知なので，本設定で自由度を設定して $t$ 分布を使う必要はない。
分散成分を未知として推定した別の問題に置き換えた場合は，指定された検定・区間の点を与えない。

\ans{4}\points{7}
\begin{itemize}
\item 境界帰無仮説における棄却域 $Z>1.96$：1点。
\item 対立仮説の下での正規分布と検出力 $\Phi(2/\sqrt{15.2/K}-1.96)$：2点。
\item 検出力条件から $K\geq15.2(1.96+0.84)^2/4=29.792$：3点。
\item 最小整数 $K=30$（各群）と明記：1点。
\end{itemize}
正しい公式への代入のみで対立分布の導出がない場合は5点まで。
両群合わせて30施設と誤記した場合は，最後の1点を与えない。

\ans{5}\points{7}
\begin{itemize}
\item 予算制約を分散に代入：2点。
\item 微分または相加相乗平均から $m_*=\sqrt{c_0\sigma_e^2/(c_1\sigma_B^2)}$：2点。
\item 最小であることの確認：1点。
\item $m=30,K=40$ と予算を満たすこと：2点。
\end{itemize}
固定した施設数で $m$ を増やすだけの議論では予算制約を反映していない。
\newpage

\qhead{5（共通問題形式）}
\ans{1}\points{8}
\begin{itemize}
\item $M=(X_1+X_2+|X_1-X_2|)/2$ 等の有効な出発点：2点。
\item 正規分布の絶対値の期待値を計算し，$\E[M]=\sigma/\sqrt\pi$：3点。
\item 独立性または2次モーメントを正しく使い，$\Var(M)=\sigma^2(1-1/\pi)$：3点。
\end{itemize}
最大値の密度を積分しても同点。結果のみなら期待値1点，分散1点の計2点まで。
和と差の無相関性から独立性を述べる場合は，同時正規性を踏まえていること。

\ans{2}\points{5}
\begin{itemize}
\item 独立性を用いて $\Prob(M\leq t)=\Phi(t/\sigma)^4$：2点。
\item 誤棄却確率 $1-0.95^4\simeq0.185$ と計算：3点。
\end{itemize}
$4\times0.05=0.20$ は上界であり，この独立な設定での正確な確率ではない。その場合は最後の項目を1点までとする。

\ans{3}\points{4}
\begin{itemize}
\item 各検定の水準 $0.05/4=0.0125$：2点。
\item 付表1の $Q(2.24)=0.0125$ から棄却限界値約2.24：2点。
\end{itemize}
片側検定なので，各裾にさらに2で割った値を割り当てない。

\ans{4}\points{5}
\begin{itemize}
\item 4つの $p$ 値が順に0.0049，0.0179，0.0250，0.4207：計2点（各0.5点）。
\item Bonferroni法で $H_1$ だけ棄却：1点。
\item Holm法で第2段階で停止し，$H_1$ だけ棄却：2点。
\end{itemize}
Holm法で $H_3$ も棄却した場合，最後の項目は0点。停止規則を無視しているため，単なる数値の丸めとは扱わない。

\ans{5}\points{8}
\begin{itemize}
\item 最初の真の帰無仮説の順位 $k$ に着目し，$m-k+1\geq m_0$ を示す：3点。
\item 誤棄却があれば $\min_{j\in\mathcal T}p_j\leq\alpha/m_0$ が必要と示す：2点。
\item 和事象の上界と $p$ 値の条件からFWER $\leq m_0(\alpha/m_0)=\alpha$：3点。
\end{itemize}
全仮説が真の場合だけの証明は3点まで。独立性を追加して証明した場合は，設問の一般性を満たさないため4点まで。
\end{document}
